\documentclass{article}
\usepackage{lmodern, amssymb,amsmath, graphicx, hyperref}
% === bibliography package ===
\usepackage{natbib}
% \usepackage[colorinlistoftodos, prependcaption]{todonotes} % to use \todo 
\usepackage[margin=1in]{geometry}
\usepackage{csquotes}
\usepackage{setspace}
\hypersetup{
    colorlinks=true,
    linkcolor=blue,
    filecolor=magenta,      
    urlcolor=cyan,
    citecolor = black
}
\urlstyle{same}  % don't use monospace font for urls

  \title{A Comprehensive Analysis of Constituent Service Requests: Legislative Behavior, Representation, and Money in Politics}
    \author{Devin Judge-Lord\thanks{Ph.D. Candidate, University of Wisconsin-Madison}\and Justin Grimmer\thanks{Professor, Stanford University} \and Eleanor Neff Powell\thanks{Booth Fowler Associate Professor, University of Wisconsin-Madison}}
    \date{\today}
    
\graphicspath{{gh-pages/Figs/}}


\begin{document}

\maketitle


%%%  4600 Character Max with Spaces
%%% Must have following headers: Overview, Intellectual Merit, Broader Impacts

%\begin{abstract}
%\noindent Representation in Congress is based, in part, on legislators' ability to assist their constituents with the federal bureaucracy.  Constituency service is central to definitions of representation as well as empirical explanations for phenomena such a the incumbency advantage. And yet, too little is known about when and how elected officials contact the bureaucracy on behalf of their constituents. We assemble a massive new database of over 150,000 Congressional requests to agencies between 2007 and 2017 obtained through over 250 FOIA requests, approaching a census of all contacts with the bureaucracy. We examine when legislators contact agencies, which legislators contact those agencies, and the purpose of the contact.  Our work provides new insights into who provides constituency service, why they provide that service, and the consequences of the agency requests for how governments make decisions. % Not only do representative vary significantly in their productivity, the top advocates for individual constituents are also the top advocates for policy change, making contact with the bureaucracy an important dimension of representation.
%\end{abstract}

%\newpage
\doublespacing

\section{Overview}

Representation in Congress is based, in part, on legislators' ability to assist their constituents with the federal bureaucracy.  Constituency service is central to definitions of representation as well as empirical explanations for phenomena such a the incumbency advantage. And yet, too little is known about when and how elected officials contact the bureaucracy on behalf of their constituents. We are seeking funding to assemble a massive new database of Congressional requests to agencies between 2007 and 2018 obtained through over 250 FOIA requests, approaching a census of all contacts with the bureaucracy. We examine when legislators contact agencies, which legislators contact those agencies, and the purpose of the contact.  Our work will provides new insights into who provides constituency service, why they provide that service, and the consequences of the agency requests for how governments make decisions. 



% Not only do representative vary significantly in their productivity, the top advocates for individual constituents are also the top advocates for policy change, making contact with the bureaucracy an important dimension of representation.

\section{Intellectual Merit}
\section{Broader Impacts}

\bibliographystyle{chicago-apsr}

\bibliography{congress2018}

%\newpage
%\listoftodos[Notes]


\end{document}






